% For compatibility with acroread version 
\pdfminorversion=4
\documentclass[10pt, mathserif, aspectratio=169, t]{beamer}
\mode<presentation>

%\usefonttheme[onlymath]{serif}
\usepackage[T1]{fontenc}
%\usepackage{tgheros}
%\usepackage{tgadventor}
%\usepackage{libertine}
%\usepackage{theinhart}
\usepackage{helvet}
%\usepackage{newtxtext} %Times font only for text

% Use my theme
\usetheme{rognes}

% Used for image placement
\def\imagetop#1{\vtop{\null\hbox{#1}}}

% Misc packages
\usepackage{graphicx}
\usepackage[export]{adjustbox}
\usepackage{stmaryrd}
\usepackage{tikz}
\usepackage{booktabs}
\usepackage{multirow}
\usetikzlibrary{arrows,shapes,backgrounds,decorations,mindmap,patterns,snakes}
\usepackage{amsmath, amssymb, amsthm, amscd}
\usepackage[english]{babel}
\usepackage{algorithm}
\usepackage{pythonhighlight}
\usepackage{epstopdf}
\usepackage{multimedia}
\usepackage{soul}

\newcommand{\R}{\mathbb{R}}
\DeclareMathOperator{\Div}{\mathrm{div}}
\DeclareMathOperator{\Curl}{\mathrm{curl}}
\DeclareMathOperator{\Grad}{\mathrm{grad}}
\DeclareMathOperator{\D}{d}
\newcommand{\foralls}{\forall \;}
\newcommand{\iinner}[2]{\llangle #1, #2 \rrangle}
\newcommand{\inner}[2]{\langle #1,  #2\rangle}
\newcommand{\totald}{\mathrm{d}}
\newcommand{\refer}[1]{\begin{flushright}{\tiny \textcolor{darkgray}{[#1]}}\end{flushright}}
\newcommand{\referinline}[1]{{\tiny \textcolor{darkgray}{[#1]}}}
\newcommand{\ddt}[1]{{#1}_t}
\newcommand{\Mi}{M_i}
\newcommand{\Me}{M_e}
\newcommand{\Iion}{I_{\rm{ion}}}
\newcommand{\dx}{\, \mathrm{d}x}
\newcommand{\ds}{\, \mathrm{d} s}
\newcommand{\triang}{\mathcal{T}}
\DeclareMathOperator{\trace}{tr}
\newcommand{\duality}[2]{\langle #1, #2 \rangle}

\newcommand{\mysection}[1]{\begin{frame} \begin{center} \vspace{3em} \textbf{#1} \end{center} \end{frame}}
\newcommand{\videosection}[3]{\begin{frame} \begin{center} \vspace{3em} \href{#2}{\textcolor{rognesred}{\textbf{#1}}} \end{center} #3 \end{frame}}
\newcommand{\sectionwithrefer}[2]{\begin{frame} \begin{center} \vspace{3em} \textbf{#1} \end{center} \refer{#2} \end{frame}}

\newtheorem{prop}[theorem]{Proposition}


\tikzstyle{na} = [baseline=-.5ex]
\tikzstyle{every picture}+=[remember picture]
\everymath{\displaystyle}

\usepackage[small, figurename=,size=scriptsize,font={color=darkgray}]{caption}



%% \usepackage{listings}
%% \usepackage[most]{tcolorbox} 	%used for writing process notes, code boxes 
%% \tcbuselibrary{listings,skins}  % skins needed for code box styles
%% \usepackage{color}
%% \usepackage{amsfonts}
%% \usepackage{fancyvrb}		%provides tab and white space preservation for 
%% 				%the verbatim environment. Needed for Python
%% 				%code inclusion

\definecolor{programcode}{gray}{0.65}

\def\formtmpX#1#2{{\vskip3pt\noindent\fboxsep=0pt{\parbox{\textwidth}{\hbox to \textwidth{\hskip3pt\vbox{\raggedright\noindent\textbf{#2\vphantom{Qy}}}\hfill}}}\vskip3pt\par
\noindent\kern0pt}}

\newenvironment{programcode}[1]{\ignorespaces\def\stmtopen##1{##1}%
\formtmpX{programcode}{\centerline{\small{#1}}}}{\noindent\textcolor{programcode}{\rule{\columnwidth}{0pt}}\par\addvspace{\baselineskip}}%



\newcommand{\terminal}[1]{
  \vspace{-1em}
  \begin{programcode}{}%{Terminal window}%
    \colorbox{blue!10}{\parbox{0.98\textwidth}{\textcolor{black}{\texttt{#1}}}}
  \end{programcode}
  \vspace{-0.5em}
}

\newcommand{\emp}[1]{\texttt{#1}}

%-------------------------------------------------------------------------------
%-------------------------------------------------------------------------------

\DeclareMathOperator{\ssum}{\textstyle \sum}
\begin{document}

%-------------------------------------------------------------------------------
\cleanpage
\begin{frame}
  \begin{columns}
    \begin{column}{0.25\textwidth}
    \centering
    \includegraphics[width=\textwidth]{graphics/simula_logo_main_RGB.pdf} \\
    \vspace{1em}
    \includegraphics[width=\textwidth]{graphics/waterscales.pdf}  \\
    \vspace{1em}
    \includegraphics[height=0.3\textwidth]{graphics/flag_yellow_low.jpg} 
    \hspace{0.2em}
    \includegraphics[width=0.3\textwidth]{graphics/logo-erc.jpg} \\
    \vspace{1em}
    \includegraphics[width=0.5\textwidth]{graphics/waterscape_logo.png} \\
    \vspace{1em}
    \includegraphics[width=\textwidth]{graphics/rcn-logo.pdf} \\
    \end{column}
    \begin{column}{0.75\textwidth}
    {
    \centering

      \vspace{4em}

      {\bf Brain modelling: \\ from magnetic resonance images \\ to finite element simulation} \\

      \bigskip
      \bigskip
      

      Marie E. Rognes \\
      \bigskip
      \bigskip

      Oslo, Norway \\
      \medskip
      Jan 15, 22, and 29 2021 \\
      }
    \end{column}
  \end{columns}
\end{frame}

\cleanpage

\begin{frame}
  \frametitle{Outline of lectures}
  \begin{center}
  \includegraphics[height=2.3cm]{/home/meg/papers/current/hikeandsurf/Book/graphics/chp1/T1-image-rot-white}
  \includegraphics[height=2.3cm]{/home/meg/papers/current/hikeandsurf/Book/graphics/chp1/ernie-volume-64}
  \includegraphics[height=2.3cm]{/home/meg/papers/current/hikeandsurf/Book/graphics/chp1/soltn-t30-crop}
  \end{center}
  \begin{description}
    \item[Lectures 1-2] Quick start: From brain MRI to FEM
    \item[Lectures 3-4] The poroelastic brain  
    \item[Lectures 5-6] Introducing diffusion tensor images (anisotropy, parcellations)
      \begin{enumerate}[I]
      \item
        Brain parcellations 
      \item
        Anisotropy and the brain
      \item
        Brain physics and chemistry at other scales
      \end{enumerate}
  \end{description}
\end{frame}

\graphicspath{{/home/meg/presentations/slidesx/}}

\mysection{I: Brain parcellations}

\input{/home/meg/presentations/slidesx/mri2fem_parcellations.tex}

\mysection{II: Anisotropy in brain characteristics}

\input{/home/meg/presentations/slidesx/mri2fem_dti_anisotropy.tex}

\mysection{III: Simulating tracer influx by diffusion in the brain}

\input{/home/meg/presentations/slidesx/mri2fem_full_simulation.tex}

\mysection{Igor, bring me the last brains!}
 
\begin{frame}
\frametitle{Want to give it a try? Natural next steps:}
\begin{enumerate}
\item
  Read Chapters 5-6 of \href{https://github.com/kent-and/mri2fem}{Mardal et al (2021)}
\item
  Try out "abby" in the mri2fem data set        
\item
  Explore and have fun!                
\end{enumerate}
\end{frame}

\mysection{Brain modelling at other scales}

\sectionwithrefer{In what sense is perivascular flow driven by arterial pulsations?}{\href{https://doi.org/10.1371/journal.pone.0244442}{Daversin-Catty et al, PLoS ONE, 2020}}

\input{/home/meg/presentations/slidesx/brain_mechanics_pvs_2020.tex}

\begin{frame}
  \begin{center}
    \vspace{3em}
    \bf Osmotic versus hydrostatic forces in brain tissue
  \end{center}

  \bigskip
  
  \refer{\href{}{Ellingsrud, Dukefoss, Enger, Halnes, Pettersen and Rognes, 
  Validation and knowledge-integration of a computational framework \\
  for ionic electrodiffusion with cortical spreading depression as a
  case study, in-prep, 2021}}

  \refer{\href{}{Ellingsrud, Boull\'e, Farrell, and Rognes, Accurate
      numerical simulation of electrodiffusion and osmotic water
      movement in brain tissue, in-prep, 2021}}

\end{frame}

\input{/home/meg/presentations/slidesx/brain_accurate_mori_2021.tex}

%% %-------------------------------------------------------------------------------


\end{document}
